\documentclass[pamm,a4paper,fleqn]{w-art}

\usepackage{times,cite,w-thm}
\usepackage[T1]{fontenc}
\usepackage[utf8]{inputenc}
\usepackage{graphicx}
\usepackage{float}
\usepackage{amsmath,amssymb,amsfonts}
\usepackage{xcolor}
\usepackage{nicefrac}
\usepackage{xspace}

\DeclareRobustCommand{\mathdefault}[1]{#1}

\newcommand{\vMLabel}{[\textbf{vM}]\xspace}
\newcommand{\Kinc}{k_{\rm inc}}
\newcommand{\Jinc}{j_{\text{inc}}}
\newcommand{\JcZero}{J_c^0}
\newcommand{\Jeff}{J_{x,c}^{\rm eff}}

\newcommand{\NumCases}{45}
\newcommand{\WstegValues}{0.4, 0.5, 0.6, 1, 2}
\newcommand{\JeffBaselineMin}{1.75}
\newcommand{\JeffBaselineMax}{1.80}
\newcommand{\MeanGcLowEffect}{-27.5\%}
\newcommand{\MeanGcHighEffect}{-10.0\%}
\newcommand{\MeanKLowEffect}{-25.8\%}
\newcommand{\MeanKHighEffect}{-14.3\%}


\begin{document}

\TitleLanguage[EN]
\title[Effective crack resistance with circular inclusions]%
{Effective crack resistance of a simplified ductile
  heterogeneous material with circular inclusions}

\author{\firstname{Alexander} \lastname{Schl\"{u}ter}\inst{1,}%
  \footnote{Corresponding author:
    \ElectronicMail{alexander.schlueter@tu-darmstadt.de}}}

\author{\firstname{Ronjit} \lastname{Medda}\inst{1}}
\author{\firstname{Ralf} \lastname{M\"{u}ller}\inst{1}}

\address[\inst{1}]{\CountryCode[DE]Institut f\"ur Mechanik,
  FB Bau- und Umweltingenieurwissenschaften, TU Darmstadt,
  Franziska-Braun-Stra{\ss}e 7, 64287 Darmstadt, Germany}

\AbstractLanguage[EN]
\begin{abstract}
  This contribution investigates the effective crack resistance
  obtained from a numerical crack-growth experiment for a simplified
  two-dimensional ductile heterogeneous material. The resolved
  microstructure contains circular inclusions. Their stiffness and
  fracture resistance are varied independently together with the
  ligament width between neighboring inclusions. Crack evolution is
  simulated with a phase field model for ductile fracture based on
  von Mises plasticity. The maximum boundary J-integral during
  enforced macroscopic crack advance is used as the operational
  measure of effective crack resistance. The results show that
  inclusion fracture resistance and stiffness are strongly coupled.
  Weak, compliant inclusions yield the lowest effective crack
  resistance, whereas stiff, tough inclusions yield the highest
  values in the investigated parameter range. Increasing inclusion
  stiffness generally raises the resistance for weak and tough
  inclusions, while the ligament-width dependence is most pronounced
  for weak inclusions. These findings demonstrate that the effective
  response is governed by the combined inclusion properties and their
  spacing rather than by the local fracture resistance alone.
\end{abstract}

\maketitle

\section{Introduction}\label{sec:intro}

The determination of fracture resistance parameters for heterogeneous
materials remains less mature than the homogenization of elastic and
plastic material responses. Classical fracture mechanics introduces
critical
quantities such as the stress-intensity factor \(K_{Ic}\) or the
critical energy release rate \(J_c\), expressed through the criterion
\begin{equation}
  J = J_c.
  \label{eq:fracture_criterion_J}
\end{equation}
For homogeneous ductile materials this type of energetic criterion
remains meaningful at least up to the onset of monotonic crack
growth, cf. Gross and Seelig~\cite{gross_fracture_2011}. However, for
heterogeneous materials it is not obvious how a single
macroscopic fracture parameter should be defined, whether it
converges under enlargement of a representative domain, and how it
depends on the underlying microstructure.

This question is particularly relevant for additively manufactured
metallic materials, where pores, inclusions, and locally different
material regions may be introduced by the process or by an
intentional multi-material design. For instance,
wire arc additive manufacturing enables the fabrication of complex
metallic structures
and also motivates fracture-mechanics-based assessment in materials
with non-trivial microstructure~\cite{hauser_porosity_2021}. A
practical macroscopic analysis then requires parameters that condense
the fracture process at the microscale into a usable crack-resistance
measure.

Different routes toward effective fracture properties have been
proposed. Schneider~\cite{schneider_fftbased_2020} and Michel and
Suquet~\cite{michel_merits_2022}, among others, discuss variational
definitions related to the surface energy of a heterogeneous medium.
In the present work we follow the numerical experiment introduced by
Hossain et al.~\cite{hossain_effective_2014}. A representative
material region is forced to separate under surfing boundary
conditions, while the macroscopic energetic driving force is measured
by a boundary J-integral. The maximum value attained during the
process is interpreted as an effective crack resistance for this
particular experiment. This terminology is used deliberately: the
value is an operational crack-resistance measure for the prescribed
microstructure and loading path, not a proof of a homogenized
fracture toughness.

Fracture in elastically heterogeneous materials has already been
addressed with phase field models. Hossain et
al.~\cite{hossain_effective_2014} demonstrated toughening and
toughness asymmetry caused by elastic heterogeneity; associated
studies examined stress-fluctuation and crack-renucleation mechanisms
in layered materials and their anisotropic effective
toughness~\cite{hsueh_stress_2018,brach_anisotropy_2019}. Kuhn and
M\"uller~\cite{kuhn_discussion_2016} investigated crack-driving
configurational forces and far-field J-integrals for heterogeneous
elastic structures with spatially varying stiffness and fracture
resistance, while Bohnen et al.~\cite{bohnen_simulation_2024} studied
crack propagation and toughness determination in heterogeneous
materials with a fracture phase field. Inclusion-containing elastic
microstructures have also been treated directly: Singh et
al.~\cite{singh_optimization_2023} optimized stiff elliptical
inclusions for enhanced resistance to brittle bulk fracture.

Related applications of the present numerical experiment to porous
materials considered both elastic and nonlinear material
descriptions~\cite{schluter_determination_2025,%
  schluter_influence_2025}.
However, to the authors' knowledge, effective crack resistance for
explicit inclusions embedded in a plastically deforming matrix has
not yet been studied by this phase field numerical experiment. Here
we consider a simplified two-dimensional material with circular
inclusions. The matrix follows a von Mises-type
ductile phase field fracture model, while the inclusions are assigned
independently varied stiffness and fracture resistance. The study
investigates how these inclusion properties and the ligament width
between
inclusions affect the computed effective crack resistance. The
manuscript is structured as follows:
Section~\ref{sec:numericalExperiment}
recalls the numerical experiment,
Section~\ref{sec:simulation_phase_field_models} summarizes the
ductile phase field model and implementation,
Section~\ref{sec:2Dresults} presents the inclusion study, and
Section~\ref{sec:conclusions} summarizes the observed trends.

\section{A Numerical Experiment to Determine Effective Crack
  Resistance}\label{sec:numericalExperiment}

\begin{figure}[H]
  \centering
  \resizebox{0.6\textwidth}{!}{\input{pics/num_experiment.pdf_t}}
  \caption{Sketch of the numerical experiment introduced by Hossain
    et al.~\cite{hossain_effective_2014}. A material region with
    explicitly resolved microstructure is separated in a simulation
    while the boundary J-integral records the macroscopic
    crack-driving force.}
  \label{fig:numerical_experiment}
\end{figure}

The numerical experiment is shown schematically in
Fig.~\ref{fig:numerical_experiment}. A region containing the
particular
microstructure is embedded in a surrounding effective medium.
Macroscopic mode-I crack growth is enforced by surfing boundary
conditions, where a fictitious crack tip travels through the domain
with velocity \(\dot{x}_{\rm bc}\). The boundary displacement field
is prescribed from the asymptotic mode-I crack-tip solution,
\begin{align}
  u_x & = \frac{K_I^{\rm bc}}{2\mu_{\rm bc}}
  \sqrt{\frac{r}{2\pi}}
  \left(\kappa-\cos\varphi\right)
  \cos\frac{\varphi}{2}, \\
  u_y & = \frac{K_I^{\rm bc}}{2\mu_{\rm bc}}
  \sqrt{\frac{r}{2\pi}}
  \left(\kappa-\cos\varphi\right)
  \sin\frac{\varphi}{2},
  \label{eq:surfing-bcs}
\end{align}
with \(\kappa=(3-\nu_{\rm bc})/(1+\nu_{\rm bc})\) and
\(\mu_{\rm bc}=E_{\rm bc}/(2(1+\nu_{\rm bc}))\). Here \(r\) and
\(\varphi\) are polar coordinates centered at a fictitious crack
tip. The model is expressed in terms of a reference fracture
resistance \(J_c^0\), a reference shear modulus \(\mu_0\), and a
reference length \(L\). The boundary condition parameters used in
this work are
\begin{equation}
  \nu_{\rm bc}=0.25,\qquad
  E_{\rm bc}=2.5\mu_0,\qquad
  K_I^{\rm bc}=\sqrt{2.5J_c^0\mu_0}.
\end{equation}
The microscopic crack path is not prescribed; it follows from the
phase field fracture simulation.

During the separation process, the macroscopic driving force is
evaluated as the \(x\)-component of the J-integral at the outer
boundary,
\begin{equation}
  J_x = \int_{\partial\Omega}
  \left(
  \psi_{\rm el}\boldsymbol{1}
  - (\nabla\boldsymbol{u})^T\boldsymbol{\sigma}
  \right)
  \boldsymbol{n}\,{\rm d}A \cdot \boldsymbol{e}_x .
  \label{eq:J-integral}
\end{equation}
The effective crack resistance used in this work is then defined as
the maximum value of \(J_x\) during the enforced crack advance,
\begin{equation}
  \Jeff = \max_t J_x(t).
  \label{eq:effective_resistance}
\end{equation}
as originally proposed by Hossain
et al.~\cite{hossain_effective_2014}. This definition captures the
largest momentary macroscopic energy flux required
to force complete separation of the resolved material region. It is
a property of the microscopic crack path, microstructure,
constitutive model, and enforced macroscopic crack-growth direction.

\section{Simulation of Fracture with Phase Field
  Models}\label{sec:simulation_phase_field_models}


\subsection{Ductile phase field fracture
  model}\label{sec:phase_field_model_ductile}

Fracture is modeled by the phase field
formulation of Noll et al.~\cite{noll_3d_2020}. Small strains are
assumed and the total strain is additively decomposed as
\begin{equation}
  \boldsymbol{\varepsilon} =
  \boldsymbol{\varepsilon}^e+\boldsymbol{\varepsilon}^p .
\end{equation}
The unfractured material follows von Mises plasticity with linear
isotropic hardening. The yield stress is
\begin{equation}
  \sigma_Y(\alpha)=\sigma_y+H\alpha,
\end{equation}
where \(\alpha\) is the accumulated plastic variable, \(H\) the
hardening modulus, and \(\sigma_y\) the initial yield stress. With
the deviatoric part of the stress tensor
\(\boldsymbol{s}=\operatorname{dev}(\boldsymbol{\sigma})\), the yield
function reads
\begin{equation}
  f(\boldsymbol{\sigma},\alpha)
  = \|\boldsymbol{s}\|
  - \sqrt{\frac{2}{3}}
  \left(\sigma_y+H\alpha\right).
\end{equation}
The evolution of the plastic variables follows the usual associative
\(J_2\) flow rule together with the Karush--Kuhn--Tucker
conditions, see e.g.~\cite{simo_computational_1998}.

Fracture is described by a scalar phase field
\(s(\boldsymbol{x},t)\in[0,1]\), where \(s=1\) denotes intact
material and \(s=0\) fully fractured material. The degraded elastic
energy density is
\begin{equation}
  \psi_{\rm el}(
  \boldsymbol{\varepsilon},
  \boldsymbol{\varepsilon}^p,
  s)
  = g(s)w_{\rm el}(
  \boldsymbol{\varepsilon},
  \boldsymbol{\varepsilon}^p),
  \qquad
  w_{\rm el}(\boldsymbol{\varepsilon},\boldsymbol{\varepsilon}^p)
  = \frac{1}{2}
  \left(\boldsymbol{\varepsilon}-\boldsymbol{\varepsilon}^p\right)
  :
  \left(
  \boldsymbol{\mathbb{C}}:
  \left(
  \boldsymbol{\varepsilon}-\boldsymbol{\varepsilon}^p
  \right)
  \right).
\end{equation}
For the plane-strain formulation, the components of the fourth-order
elasticity tensor are
\begin{equation}
  \mathbb{C}_{ijkl}
  =(\lambda+\mu)\delta_{ij}\delta_{kl}
  +2\mu\left[
  \frac{1}{2}
  \left(\delta_{ik}\delta_{jl}+\delta_{il}\delta_{jk}\right)
  -\frac{1}{2}\delta_{ij}\delta_{kl}
  \right],
  \qquad i,j,k,l\in\{1,2\},
  \label{eq:plane_strain_elasticity_tensor}
\end{equation}
where \(\delta_{ij}\) is the Kronecker delta.
The \(J_2\) plastic strain is deviatoric, i.e.,
\(\operatorname{tr}(\boldsymbol{\varepsilon}^p)=0\).
The cubic degradation function is
\begin{equation}
  g(s)=\beta(s^3-s^2)+3s^2-2s^3+\eta .
  \label{eq:deg_grad}
\end{equation}
In this work \(\beta=0.2\) and \(\eta=10^{-5}\). The plastic
contribution is degraded in the same manner as the strain energy,
\begin{equation}
  \psi_{\rm pl}(\alpha,s)=g(s)w_{\rm pl}(\alpha),
  \qquad
  w_{\rm pl}(\alpha)=\frac{1}{2}H\alpha^2+\sigma_y\alpha .
\end{equation}
The regularized fracture energy density is
\begin{equation}
  \psi_{\rm fr}
  =J_c\left(
  \frac{(1-s)^2}{2\epsilon}
  +2\epsilon\|\nabla s\|^2
  \right),
\end{equation}
where \(J_c\) is the local fracture resistance and \(\epsilon\)
controls the width of the diffuse crack zone. The resulting phase
field evolution equation is
\begin{align}
  \dot{s}=-M\left[
    g'(s)\left(w_{\rm el}+w_{\rm pl}\right)
    -J_c\left(2\epsilon\Delta s+\frac{1-s}{2\epsilon}\right)
    \right],
  \label{eq:evolution_s}
\end{align}
with homogeneous Neumann boundary conditions for \(s\). Quasi-static
mechanical equilibrium is imposed by
\begin{equation}
  \operatorname{div}\boldsymbol{\sigma}=\boldsymbol{0}.
  \label{eq:equilibrium}
\end{equation}
For the two-dimensional plane-strain formulation, the stress is
computed as
\begin{equation}
  \boldsymbol{\sigma}=g(s)\left[
    (\lambda+\mu)
    \operatorname{tr}(\boldsymbol{\varepsilon})\boldsymbol{1}
    +2\mu\left(
    \operatorname{dev}(\boldsymbol{\varepsilon})
    -\boldsymbol{\varepsilon}^p
    \right)
    \right].
  \label{eq:sigma_plast}
\end{equation}
\subsection{Numerical
  implementation}\label{sec:numerical_implementation}

The weak forms of the phase field evolution equation and the
equilibrium equation are implemented in FEniCSx~\cite{_fenicsx_}. The
displacement, phase field, and plastic internal variables are solved
monolithically in pseudo-time. The phase field equation is
discretized by a backward Euler scheme, and the nonlinear system is
solved by Newton's method. We employ automatic time stepping that
chooses the time step such that the number of Newton iterations
remains within
the range of four to eight.

First-order Lagrange finite elements are used on triangular meshes.
The mesh is refined in the expected fracture region such that the
phase field length scale $\epsilon$ is resolved. In
Fig.~\ref{fig:setup_inclusions}~b), the minimum edge length is
\(1.30\times 10^{-2}L\), the mean edge length is \(3.24\times
10^{-2}L\), and \(\epsilon=0.1L\).

\section{Effective Crack Resistance of a Simplified Two-Dimensional
  Heterogeneous Material}\label{sec:2Dresults}

This section applies the numerical experiment to a material with
circular inclusions and examines how ligament width, inclusion
stiffness, and inclusion fracture resistance affect the effective
crack resistance.
All evaluated crack patterns are qualitatively similar to
Fig.~\ref{fig:setup_inclusions}~b). The crack propagates
horizontally and separates the central row of inclusions from the
surrounding material.

\subsection{Geometry and material parameters}\label{sec:geometry}

\begin{figure}[H]
  \centering
  \begin{minipage}[t]{0.49\textwidth}
    \centering
    \textbf{a)}\\[-0.3em]

    \makebox[\textwidth][c]{%
      \includegraphics[
        height=0.24\textheight,
        trim=115 5 115 5,
        clip]%
      {pics/material_layout_inclusions.png}}
  \end{minipage}
  \hfill
  \begin{minipage}[t]{0.49\textwidth}
    \centering
    \textbf{b)}\\[-0.3em]

    \makebox[\textwidth][c]{%
      \includegraphics[
        height=0.24\textheight,
        trim=150 5 30 5,
        clip]%
      {pics/failure_pattern_phasefield.png}}
  \end{minipage}
  \caption{Representative setup for the inclusion study. a) Full
    computational domain, including the resolved material region with
    circular inclusions and the surrounding effective medium. b) Final
    phase field pattern and triangular finite-element mesh for one
    stored simulation, plotted with a blue-to-red scale for the phase
    field \(s\). The shown simulation has \(w_s=1.0L\),
    \(\Kinc=0.5\), and \(\Jinc=0.5\).}
  \label{fig:setup_inclusions}
\end{figure}

The investigated two-dimensional material consists of square unit
cells with a centered circular inclusion of diameter \(L\). The
ligament width between neighboring inclusions is denoted by \(w_s\),
so that the cell size is \(w_c=L+w_s\). A \(6\times3\) array of unit
cells forms the resolved region of the numerical experiment; see
Fig.~\ref{fig:setup_inclusions}~a). This resolved region is embedded
in a surrounding effective medium, for which the effective elastic
properties are determined by numerical homogenization.
Plane strain conditions are assumed.

An initial horizontal crack is placed at \(y=0\), at the height of
the center row of inclusions, and extends to the left boundary of the
resolved region.

The simulated crack path remains horizontal and separates the central
row of inclusions in all cases considered here; see
Fig.~\ref{fig:setup_inclusions}~b).


The matrix properties are set equal to reference values \(\mu_0\) and
\(J_c^0\),
\begin{equation}
  \lambda_{\rm mat}=\mu_0,\qquad
  \mu_{\rm mat}=\mu_0,\qquad
  J_{c,\rm mat}=J_c^0.
\end{equation}
The yield stress \(\sigma_y\) and hardening modulus \(H\) are the same
in the matrix, inclusions, and surrounding effective medium, i.e.,
\begin{equation}
    \sigma_y=\sqrt{\frac{J_c^0\mu_0}{L}},\qquad
  H=0.22\mu_0.
\end{equation}
The phase field mobility is
\begin{equation}
  M=1000\frac{\dot{x}_{\rm bc}}{J_c^0}.
\end{equation}
The parameter study conducted in this work varies the ligament width
\begin{equation}
  w_s/L \in \{\WstegValues\},
\end{equation}
the inclusion stiffness factor
\begin{equation}
  \Kinc \in \{0.5,1.0,1.5\},
\end{equation}
and the inclusion fracture resistance factor
\begin{equation}
  \Jinc \in \{0.5,1.0,1.5\}.
\end{equation}
Thus, the inclusion Lam\'e parameters are obtained by scaling the
matrix values as
\begin{equation}
  \lambda_{\rm inc}=\Kinc\lambda_{\rm mat}, \qquad
  \mu_{\rm inc}=\Kinc\mu_{\rm mat},
\end{equation}
and the inclusion fracture resistance is
\begin{equation}
  J_{c,\rm inc}=\Jinc J_{c,\rm mat}.
\end{equation}
For each combination of wall width \(w_s\) and stiffness contrast
\(\Kinc\), a preliminary
effective-stiffness computation provides the Lam\'e parameters of the
surrounding homogeneous material as mentioned above.
% The cases with
% \(w_s=0.75L\) were excluded from the evaluation because this data
% point showed problematic behavior in the local result snapshot;
% datasets with \(w_s>2.0L\) were also omitted to focus on the
% interaction-dominated ligament-width range.

\subsection{Crack tracking and local J-integral response}

At each accepted time step, the phase field is interpolated to the
mesh nodes. A node is classified as fully damaged when
\(|s|\le 0.005\). The microscopic crack-tip position is then defined
as the largest \(x\)-coordinate among all fully damaged nodes,
\begin{equation}
  x_{\rm ct}(t)
  =\max\left\{
  x_i:\left|s(\boldsymbol{x}_i,t)\right|\le0.005
  \right\}.
  \label{eq:crack_tip_position}
\end{equation}
For a distributed mesh, the maximum is taken globally over all MPI
processes.

\begin{figure}[H]
  \centering

  \includegraphics[width=0.55\textwidth]%
  {pics/fig09_crack_tip_tracking.png}
  \caption{Position of the surfing boundary condition
    \(x_{\rm bc}\) and extracted crack-tip position
    \(x_{\rm ct}\) for the stiff, weak inclusion case
    \(w_s=1.0L\), \(\Kinc=1.5\), and \(\Jinc=0.5\).}
  \label{fig:crack_tip_tracking}
\end{figure}

Figure~\ref{fig:crack_tip_tracking} shows that the measured crack tip
follows the moving boundary-condition field \(x_{\rm bc}\). The
microscopic crack-tip motion is unsteady because the crack interacts
with the inclusions.

\begin{figure}[H]
  \centering

  \includegraphics[width=0.72\textwidth]%
  {pics/fig10_Jx_xct_selected_cases.png}
  \caption{J-integral over crack-tip position for \(w_s=1.0L\): the
    reference material contrast
    \((\Kinc,\Jinc)=(1.0,1.0)\), stiff, weak inclusions
    \((1.5,0.5)\), and compliant, tough inclusions with
    increased fracture resistance \((0.5,1.5)\).}
  \label{fig:Jx_xct_stiffness}
\end{figure}

For fixed \(w_s=1.0L\), Fig.~\ref{fig:Jx_xct_stiffness} contrasts the
recorded J-integral \(J_x\) response for the reference case with
homogeneous material, i.e., \(\Kinc=1.0\), \(\Jinc=1.0\),
with two distinct modifications: a stiff, weak inclusion phase with
\(\Kinc=1.5\), \(\Jinc=0.5\), and a compliant, tough inclusion phase
with \(\Kinc=0.5\), \(\Jinc=1.5\).
For the reference case, the J-integral response is almost constant
during crack propagation at a value of
\(J_{x,c,\rm ref}^{\rm eff}\approx1.44J_c^0\),
which is expected for a homogeneous material.
Both the stiff, weak inclusions and the compliant, tough inclusions
show a typical periodic pattern of the J-integral response that is
characteristic for crack growth in the given material.
In both cases, the peak values of the J-integral defining the
effective crack resistance according to
\eqref{eq:effective_resistance} are higher than the reference value.
Interestingly, the graphs for the two cases show a phase shift, with
peaks for the compliant, tough inclusions occurring when the crack
tip enters the inclusion, i.e., for
\[
  x_{\rm ct}=2.5L+k(2.0L),\qquad k\in\{0,1,2,\ldots\},
\]
while the J-integral for the stiff, weak inclusions is minimal at
these positions. In the latter case, the J-integral peaks occur just
before the crack tip exits the inclusion into the more compliant
matrix material, i.e., for
\[
  x_{\rm ct}=3.5L+k(2.0L),\qquad k\in\{0,1,2,\ldots\}.
\]
This has also been observed and
discussed in \cite{kuhn_discussion_2016} for elastic materials.



% \begin{figure}[H]
%   \centering

%   \includegraphics[width=0.9\textwidth]%
%   {pics/fig11_Jx_sections_Jcinc_0p5.png}
%   \caption{Section of the J-integral \(J_x\) versus crack-tip
%     position \(x_{\rm ct}\) for \(\Jinc=0.5\). The evaluated
%     interval
%     extends from the ligament midpoint before inclusion 1 to the
%     ligament midpoint before inclusion 3. Thus, it contains the
%     interactions with inclusions 1 and 2. For \(w_s=L\), this
%     corresponds to \(2\le x_{\rm ct}/L\le6\). The three
%     panels show the available simulations for \(\Kinc=0.5\),
%     \(1.0\), and
%     \(1.5\); within each panel the curves correspond to different
%     ligament widths \(w_s\).}
%   \label{fig:Jx_sections_Jcinc_0p5}
% \end{figure}



% \begin{figure}[H]
%   \centering

%   \includegraphics[width=0.9\textwidth]%
%   {pics/fig13_Jx_sections_Jcinc_1p5.png}
%   \caption{Section of the J-integral \(J_x\) versus crack-tip
%     position \(x_{\rm ct}\) for \(\Jinc=1.5\). The evaluated
%     interval
%     extends from the ligament midpoint before inclusion 1 to the
%     ligament midpoint before inclusion 3. Thus, it contains the
%     interactions with inclusions 1 and 2. For \(w_s=L\), this
%     corresponds to \(2\le x_{\rm ct}/L\le6\). The three
%     panels show the available simulations for \(\Kinc=0.5\),
%     \(1.0\), and
%     \(1.5\); within each panel the curves correspond to different
%     ligament widths \(w_s\). The vertical axis is restricted to
%     \(0\le J_x/\JcZero\le3.2\).}
%   \label{fig:Jx_sections_Jcinc_1p5}
% \end{figure}

% Figures~\ref{fig:Jx_sections_Jcinc_0p5}--
% \ref{fig:Jx_sections_Jcinc_1p5}
% show only the characteristic J-integral section between inclusions
% one and three along the crack path for all parameter combinations.

% In Fig.~\ref{fig:Jx_sections_Jcinc_0p5}, which shows weak
% inclusions \((\Jinc=0.5)\), it can be observed that
% smaller ligament widths \(w_s\) lead to lower peaks
% in the J-integral record, i.e. lower effective crack resistance.
% The influence of the stiffness contrast \(\Kinc\) is less
% pronounced but still visible, with stiff inclusions showing higher
% peaks than compliant inclusions.



% The evaluated section contains the interactions
% with inclusions 1 and 2 and ends at the ligament midpoint before
% inclusion 3. The figures make the local crack-inclusion interactions
% visible: the
% maxima used in Eq.~\eqref{eq:effective_resistance} are not single
% smooth material responses, but arise from trapping, deflection, and
% re-acceleration of the crack as it crosses the inclusion array. This
% section lies beyond the initial transient while retaining repeated
% crack--inclusion interactions. Larger \(w_s\) separates individual
% interaction events more clearly.

\subsection{Effective crack resistance}

\begin{figure}[H]
  \centering

  \includegraphics[width=0.95\textwidth]%
  {pics/fig15_Jeff_wsteg_all_Jcinc.png}
  \caption{Effective crack resistance normalized by the reference
    value \(J_{x,c,\rm ref}^{\rm eff}=1.44J_c^0\). Each panel fixes
    the inclusion fracture-resistance factor \(\Jinc\) and compares
    the three inclusion stiffness factors.}
  \label{fig:Jeff_wsteg_families}
\end{figure}

Figure~\ref{fig:Jeff_wsteg_families} summarizes the effective crack
resistance determined from~\eqref{eq:effective_resistance} for each
parameter combination.

The effective crack resistance of the reference case
\((\Kinc,\Jinc)=(1.0,1.0)\) is denoted by
\(J_{x,c,\rm ref}^{\rm eff}\). It defines the normalization in
Fig.~\ref{fig:Jeff_wsteg_families} and is chosen as
\begin{equation}
  J_{x,c,\rm ref}^{\rm eff}=1.44J_c^0.
\end{equation}

For weak inclusions \((\Jinc=0.5)\), the effective crack resistance
\(\Jeff\) generally increases with ligament width \(w_s\) and
approaches a plateau; see
Fig.~\ref{fig:Jeff_wsteg_families}~a). Stiffer inclusions reach this
regime at smaller \(w_s\) than compliant inclusions. The stiffest
inclusions give the highest effective resistance at every investigated
ligament width. At \(w_s=2L\), the effective resistances for
\(\Kinc=0.5\), \(1.0\), and \(1.5\) seem to converge to approximately
\(1.00J_{x,c,\rm ref}^{\rm eff}\), respectively.

When the inclusion and matrix fracture resistances are equal
\((\Jinc=1.0)\), the effective crack resistance depends only weakly
on the ligament width for all three stiffness contrasts; see
Fig.~\ref{fig:Jeff_wsteg_families}~b). This occurs in spite of the
fact that the J-integral response shows a clear
periodic pattern of peaks and valleys for all ligament widths, as
illustrated in Fig.~\ref{fig:Jx_sections_Jcinc_1p0}. Compliant
inclusions yield a lower effective resistance 
of 
\(1.00J_{x,c,\rm ref}^{\rm eff}\) to
\(1.07J_{x,c,\rm ref}^{\rm eff}\), while stiff inclusions produce
values between \(1.13J_{x,c,\rm ref}^{\rm eff}\) and
\(1.15J_{x,c,\rm ref}^{\rm eff}\).

\begin{figure}[H]
  \centering

  \includegraphics[width=0.9\textwidth]%
  {pics/fig12_Jx_sections_Jcinc_1p0.png}
  \caption{Section of the J-integral \(J_x\) versus crack-tip
    position \(x_{\rm ct}\) for \(\Jinc=1.0\). The evaluated
    interval
    extends from the ligament midpoint before inclusion 1 to the
    ligament midpoint before inclusion 3. Thus, it contains the
    interactions with inclusions 1 and 2. For \(w_s=L\), this
    corresponds to \(2\le x_{\rm ct}/L\le6\). The
    three panels show the available simulations for \(\Kinc=0.5\),
    \(1.0\), and
    \(1.5\); within each panel the curves correspond to different
    ligament widths \(w_s\).}
  \label{fig:Jx_sections_Jcinc_1p0}
\end{figure}

The tough-inclusion cases \((\Jinc=1.5)\) show a more complex
ligament-width dependence; see
Fig.~\ref{fig:Jeff_wsteg_families}~c). For compliant inclusions
\((\Kinc=0.5)\), the effective resistance increases from
\(1.50J_{x,c,\rm ref}^{\rm eff}\) at \(w_s=0.4L\) to
\(1.58J_{x,c,\rm ref}^{\rm eff}\) at \(w_s=2L\). For
matrix-matched stiffness \((\Kinc=1.0)\), it remains nearly constant
between \(1.75J_{x,c,\rm ref}^{\rm eff}\) and
\(1.78J_{x,c,\rm ref}^{\rm eff}\). For stiff inclusions
\((\Kinc=1.5)\), the resistance is highest at the smallest ligament
width, where it reaches
\(2.10J_{x,c,\rm ref}^{\rm eff}\), and decreases to
\(1.94J_{x,c,\rm ref}^{\rm eff}\) at \(w_s=2L\). The stiff
inclusions therefore give the highest effective resistance, and the
compliant inclusions the lowest, at every investigated ligament
width.


\begin{figure}[H]
  \centering

  \includegraphics[width=0.95\textwidth]%
  {pics/fig14_combined_parameter_contribution.png}
  \caption{Normalized effective crack resistance at selected
    ligament widths: a) \(w_s=0.5L\), b) \(w_s=1.0L\), and c)
    \(w_s=2.0L\). Each cell shows \(\Jeff\) divided by the
    fixed reference value
    \(J_{x,c,\rm ref}^{\rm eff}=1.44J_c^0\).
    Weak, compliant inclusions \((0.5,0.5)\) give the lowest
    response, whereas stiff, tough inclusions \((1.5,1.5)\) give the
    largest response in all three panels.}
  \label{fig:combined_parameter_contribution}
\end{figure}

Figure~\ref{fig:combined_parameter_contribution} presents the same
results for three selected ligament widths. Each cell contains
\(\Jeff\) divided by the fixed reference value
\(J_{x,c,\rm ref}^{\rm eff}=1.44J_c^0\). This normalization
isolates the relative effects of inclusion stiffness and fracture
resistance.

In all three panels, the weak, compliant combination
\((\Kinc,\Jinc)=(0.5,0.5)\) yields the lowest effective crack
resistance, while the stiff, tough combination
\((\Kinc,\Jinc)=(1.5,1.5)\) yields the highest. For weak and tough
inclusions, increasing \(\Kinc\) increases the resistance. When
\(\Jinc=1.0\), the reference stiffness gives the lowest response,
although the variation with stiffness remains much smaller than for
tough inclusions. Increasing \(\Jinc\) to \(1.5\) produces the
largest gain for every stiffness, whereas the difference between
\(\Jinc=0.5\) and \(1.0\) is comparatively small. Even without an
increase in inclusion fracture resistance, stiff inclusions raise
the effective resistance by approximately \(15\%\) at \(w_s=0.5L\).


\section{Conclusions}\label{sec:conclusions}

This work applied a numerical crack-growth experiment to determine
the effective crack resistance of a simplified two-dimensional
ductile material with circular inclusions. Inclusion stiffness,
inclusion fracture resistance, and ligament width were varied
independently. The effective crack resistance was defined by the
maximum boundary J-integral attained while the crack was driven
through the resolved microstructure by a surfing boundary condition.

For every investigated parameter combination, the crack propagated
horizontally along the central row of inclusions. Its interaction
with the periodic microstructure produced a J-integral
response characteristic of the microstructure and parameter set. The
weak, compliant combination produced the lowest effective crack
resistance, whereas the stiff, tough combination produced the highest
values throughout the selected ligament-width range.

The influence of ligament width depended on the inclusion
properties. For weak inclusions, the effective resistance generally
increased with ligament width and approached a saturation value.
Stiffer weak inclusions reached this regime at smaller ligament
widths and gave a larger resistance than compliant weak inclusions.
When the inclusion and matrix fracture resistances were equal, the
effective resistance was nearly independent of ligament width for
all three stiffness contrasts. Nevertheless, stiff inclusions still
increased the resistance by up to \(15\%\) relative to the reference
case. For tough inclusions, compliant and matrix-matched cases varied
only moderately with ligament width, while stiff, tough inclusions
showed the largest resistance at the smallest ligament width and a
decreasing response as the ligament width increased.

These results demonstrate that the effective crack resistance is
governed by the combined inclusion properties and their spacing,
rather than by the local fracture resistance alone. The numerical
experiment therefore provides a compact means of comparing
microstructural design choices in ductile heterogeneous materials
with inclusions.

Future work should extend the study to other inclusion shapes and
pore geometries, as well as different spatial arrangements and volume
fractions. A more detailed investigation of the local stress, plastic
strain, phase field, and J-integral distributions is also needed to
explain why the individual combinations of stiffness, fracture
resistance, and ligament width produce the trends observed here. Such
analyses could clarify the roles of crack trapping, stress
redistribution, plastic dissipation, and crack-path selection, and
thereby support the development of mechanistic design rules for
heterogeneous ductile materials.

\begin{acknowledgement}

Funded by the Deutsche Forschungsgemeinschaft (DFG, German Research
Foundation) -- Project-ID 511263698 -- TRR 375.\\
Funded by the Deutsche Forschungsgemeinschaft (DFG, German Research
Foundation) -- Project-ID 523648868 -- FOR 5701.\\
Calculations for this research were conducted on the Lichtenberg high
performance computer of the TU Darmstadt.\\
This research was enabled by the FEniCSx library~\cite{_fenicsx_},
and we thank its developers and the community for their valuable
contributions.
\end{acknowledgement}

\bibliographystyle{pamm}
\bibliography{bib/literatur}

\end{document}
